\section{Derivations}
Notation: 
\begin{itemize}
    \item $\bm{\beta}_t$: the state variable
    \item $\bm{Q}_t=\bm{\Sigma}_t$ is the covariance of the state process noise. 
    \item $\bm{P}_{t|t}$ is the covariance of our estimate of $\bm{\beta}_t$.
    \item $R_{t}=\tau_t^2$ is the observation noise.
\end{itemize}
Prediction:
\begin{align*}
    \bm{\beta}_{t|t-1}&=\bm{\beta}_{t-1|t-1}\\
    \bm{P}_{t|t-1}&=\bm{P}_{t-1|t-1}+\bm{\Sigma}_{t}
\end{align*}
Update:
\begin{align*}
    \tilde{y}_t&=y_t-\bm{\beta}_{t|t-1}^T\bm{x}_t\\
    S_t&=\bm{x}_t^T\bm{\Sigma}_t\bm{x}_t+\tau_t^2\sim \mathrm{Cov}(\tilde{y}_t)\\
    K_t&=\cdots
\end{align*}
\subsection{Initialisation}
TO ADD
\subsection{PCA-KF}
The idea is that if $p$ is large, then we apply dimension reduction to $\bm{x}_t$. This assumes the data generating process doesn't change `too much' in the training set.

Observe that if $UU^T\approx I_p$, then
\[\bm{y}_t\approx\bm{\beta}_t^TUU^T\bm{x}_t+\bm{z}_t\]
This is the case, for example, if $U$ represent PCA loadings with sufficient proportion of variance retained. 
